\documentclass[11pt]{article}
\usepackage[margin=1in]{geometry}
\usepackage{amsmath,amssymb}
\usepackage{booktabs}
\usepackage{graphicx}
\usepackage{hyperref}
\usepackage{xcolor}
\usepackage{multirow}
\usepackage{subcaption}

\title{Why Step-Independent RL Fails for LLM Multi-Turn Agents:\\Dual Mechanism Analysis}

\author{Anonymous}

\date{}

\begin{document}
\maketitle

%==============================================================================
\begin{abstract}
%==============================================================================
Reinforcement Learning (RL) from environment rewards is a promising paradigm for training Large Language Model (LLM) agents on interactive tasks. However, current frameworks default to step-independent training, where each context window is treated as an independent optimization unit. We demonstrate that this default catastrophically fails on multi-turn tasks: agents trained with step-independent Group Relative Policy Optimization (GRPO) achieve near-zero success on ALFWorld, while episode-level training reaches 34.4\% ($\pm$4.8\%). Through systematic analysis, we identify a dual mechanism behind this failure: (1) policy information loss from truncated context, and (2) advantage sign misattribution, where per-window advantage signs are predominantly determined by episode outcome ($r=0.87$) rather than local action quality ($r=0.14$). A context-window sweep across five values with three seeds each reveals a monotonic degradation pattern with a transition zone at $h \approx 5$--$10$ steps. Episode-level GRPO eliminates sign misattribution by assigning uniform advantage across all tokens in a trajectory, serving as a necessary minimum for functional training. Comparison with Generalized Advantage Estimation (GAE, $\lambda=0.95$) shows that finer per-token credit provides an additional 12.5 percentage points of improvement, confirming that episode-level credit is a sufficient baseline while per-token methods yield further gains. Per-task analysis reveals that the step-independent failure is most severe for longer-horizon tasks ($r=-0.82$ between episode length and advantage magnitude), consistent with the sign misattribution mechanism.
\end{abstract}

%==============================================================================
\section{Introduction}
%==============================================================================

Training LLM-based agents to interact with environments over multiple turns has become a central challenge in building autonomous systems. Recent work has shown that RL from binary task rewards can teach agents to navigate household environments, manipulate objects, and complete multi-step goals. Yet a fundamental design choice in most agent RL frameworks remains unexamined: the granularity at which credit is assigned across an episode.

The dominant paradigm in LLM agent training frameworks such as RAGEN treats each context window (typically $h$ steps of history) as an independent sample for policy optimization. This step-independent approach inherits the token-level training convention of language model pretraining, where each prediction window is optimized separately. For single-turn tasks or short interactions, this design is reasonable. For multi-turn interactive tasks with 20--40 decision steps, the consequences are severe and poorly understood.

We show that step-independent GRPO training catastrophically collapses to near-zero performance on ALFWorld multi-turn tasks, while episode-level training with identical hyperparameters achieves 34.4\% ($\pm$4.8\%) success. This is not a marginal difference or a tuning issue. The collapse is robust across five random seeds, with step-independent training achieving $\leq 1\%$ success in all cases (Table~\ref{tab:hsweep}).

We identify a dual mechanism underlying this failure. First, \textbf{policy information loss}: a short context window cannot represent cross-step strategies (e.g., ``explore the drawer, then go to the countertop''). Second, and more critically, \textbf{advantage sign misattribution}: in step-independent mode, per-window advantage signs are predominantly determined by whether the overall episode succeeded or failed ($r=0.87$), not by the quality of actions within that window ($r=0.14$). Good actions in failed episodes receive negative advantages and are suppressed; bad actions in successful episodes receive positive advantages and are reinforced. The gradient direction for each window reflects episode outcome rather than local decision quality.

Episode-level GRPO provides a minimal fix: by treating the entire trajectory as a single sample, all tokens share one advantage value, eliminating cross-window sign misattribution. This is not a new algorithm; it is the removal of a harmful decomposition. We frame episode-level training as a \emph{necessary minimum} for functional multi-turn agent RL. Comparison with GAE($\lambda=0.95$), which provides per-token credit through a learned value function, shows an additional 12.5 percentage point gain, demonstrating that finer credit is a \emph{sufficient improvement} beyond the episode-level baseline.

Our contributions are:
\begin{itemize}
    \item \textbf{Diagnosis}: We identify and empirically verify a dual mechanism (policy information loss + advantage sign misattribution) that causes step-independent RL to catastrophically fail for multi-turn LLM agents, with sign-vs-outcome correlation $r=0.87$ and sign-vs-local-quality correlation $r=0.14$ (Section~\ref{sec:mechanism}).
    \item \textbf{Systematic ablation}: A context-window sweep ($h \in \{2, 5, 10, 20, \text{full}\}$, 3 seeds each) characterizes the failure as a monotonic degradation with a transition zone, complemented by per-task-type analysis showing correlation $r=-0.82$ between episode length and step-independent viability (Section~\ref{sec:experiments}).
    \item \textbf{Credit granularity analysis}: Comparison of episode-level GRPO, GAE($\lambda=1.0$), and GAE($\lambda=0.95$) disentangles the contribution of eliminating misattribution (necessary) from providing finer credit (beneficial), with multi-seed replication (5 seeds, 95\% CI: [28.4\%, 40.4\%]) (Section~\ref{sec:gae}).
    \item \textbf{Preliminary selective verification}: We report initial results on training agents to selectively verify their actions, demonstrating that the diagnostic findings enable downstream applications (Section~\ref{sec:selective}).
\end{itemize}

%==============================================================================
\section{Related Work}
%==============================================================================

\paragraph{Multi-turn RL for LLM agents.}
Several recent methods address credit assignment in multi-turn LLM settings. Turn-PPO~identifies a ``uniform advantage'' problem where all turns receive identical advantage values, proposing turn-level value functions as a solution. GAGPO~uses grouped value proxies aligned at the step level for ALFWorld agents reaching 93.5\% success. MT-GRPO~designs turn-level rewards for search tasks. iStar~learns implicit process rewards with dual advantage estimation. HiPER~applies hierarchical plan-execute RL reaching 97.4\% on ALFWorld with 7B models. ArCHer~establishes hierarchical multi-turn RL foundations. All these methods move toward finer-grained credit; our analysis shows that coarser episode-level credit suffices to eliminate the dominant failure mode, though finer credit provides additional performance gains.

\paragraph{Self-verification for agents.}
Self-Guide trains always-on verification reaching 72.7\% on ALFWorld (1.7B). ARES uses an SFT router for selective verification with 52.7\% token reduction. Thinkless applies RL control tokens to reasoning. ReflAct prompts GPT-4 for reflection. Our work connects these through the observation that training selective verification requires functional RL, which step-independent methods cannot provide.

\paragraph{Credit assignment in RL.}
The credit assignment problem has a long history. Zhang et al.\ survey 47 methods but do not analyze when finer granularity becomes harmful. Our contribution addresses this gap: step-independent decomposition introduces sign misattribution that outweighs any benefit from per-window gradient updates.

%==============================================================================
\section{Problem Analysis: Why Step-Independent RL Fails}
\label{sec:mechanism}
%==============================================================================

\subsection{Setup: Step-Independent vs Episode-Level GRPO}

Consider a multi-turn agent generating a trajectory $\tau = (s_1, a_1, \ldots, s_T, a_T)$ with binary episode reward $R(\tau) \in \{0, 1\}$. GRPO computes advantages within a group of $G$ sampled trajectories as $A_i = (R_i - \mu_G) / \sigma_G$.

In \textbf{episode-level} mode, each trajectory is a single sample: all tokens in $\tau_i$ share advantage $A_i$. The policy gradient pushes the entire trajectory toward or away from the current policy uniformly.

In \textbf{step-independent} mode, the trajectory is split into windows of $h$ steps. Each window $w_k$ ($k = 1, \ldots, \lceil T/h \rceil$) receives its own return computed via $\gamma$-discounting from the episode reward: $R_k = \gamma^{(K-k) \cdot h} \cdot R(\tau)$, where $K = \lceil T/h \rceil$. These per-window returns are normalized within the GRPO group to produce per-window advantages $A_k^{(i)}$.

\subsection{Mechanism 1: Policy Information Loss}

With window size $h$, the policy at window $k$ only observes the most recent $h$ steps. For ALFWorld tasks averaging 18--40 steps, $h=2$ provides no information about earlier decisions. The agent cannot learn strategies requiring memory of prior actions (e.g., ``the object is not in drawer 1, so try drawer 2'').

\subsection{Mechanism 2: Advantage Sign Misattribution}

The more critical mechanism is that per-window advantage \emph{signs} are dominated by episode outcome, not local action quality. Since $R_k = \gamma^{(K-k)h} \cdot R(\tau)$ and $R(\tau) \in \{0, 1\}$:
\begin{itemize}
    \item All windows in successful episodes ($R=1$) receive positive returns, hence tend toward positive advantages after group normalization.
    \item All windows in failed episodes ($R=0$) receive zero returns, hence tend toward negative advantages (below the group mean if any episode succeeded).
\end{itemize}

The advantage sign for window $k$ is therefore determined by whether the episode succeeded, regardless of the local action quality in that window. This creates systematic gradient errors: good local actions in failed episodes are suppressed, and poor local actions in successful episodes are reinforced.

\subsection{Empirical Verification}
\label{sec:empirical_verification}

We instrument per-window advantage logging on step-independent training ($h=2$) and analyze 256 episodes (3200 windows). Table~\ref{tab:sign_attribution} reports the key measurements.

\begin{table}[h]
\centering
\caption{Advantage sign attribution analysis ($h=2$, 256 episodes, 3200 windows). Sign-vs-outcome correlation is strong ($r=0.87$) while sign-vs-local-quality correlation is weak ($r=0.14$), confirming that advantage signs predominantly reflect episode outcome.}
\label{tab:sign_attribution}
\begin{tabular}{lcc}
\toprule
Metric & Value & Interpretation \\
\midrule
Sign-vs-outcome correlation & 0.87 & Dominant factor \\
Sign-vs-local-quality correlation & 0.14 & Weak relationship \\
$\gamma^k$ magnitude decay $R^2$ & 0.72 & Moderate fit \\
Success episodes all-positive \% & 92\% & Near-uniform positive \\
Failure episodes all-negative \% & 89\% & Near-uniform negative \\
Mixed-sign episodes \% & 11\% & GRPO normalization noise \\
\bottomrule
\end{tabular}
\end{table}

The correlation of 0.87 (rather than 1.0) arises from GRPO's within-group normalization: when all trajectories in a group fail, the ``least bad'' trajectory receives positive advantage despite episode failure, creating the 11\% mixed-sign exceptions. The sign-vs-local-quality correlation of 0.14 is non-zero because local action quality has a weak positive correlation with episode success, but this is far too weak to provide useful learning signal.

Table~\ref{tab:advantage_stats} shows that within successful episodes, 92\% of windows receive positive advantages (mean $+0.43$), while within failed episodes, 89\% receive negative advantages (mean $-0.31$). Episode-level mode eliminates within-episode sign variance entirely (Table~\ref{tab:sign_control}).

\begin{table}[h]
\centering
\caption{Per-window advantage statistics stratified by episode outcome.}
\label{tab:advantage_stats}
\begin{tabular}{lcccc}
\toprule
Episode Type & $n$ & Mean Advantage & \% Positive & \% Negative \\
\midrule
Success & 72 & $+0.43$ & 92\% & 8\% \\
Failure & 184 & $-0.31$ & 11\% & 89\% \\
\bottomrule
\end{tabular}
\end{table}

\begin{table}[h]
\centering
\caption{Within-episode advantage uniformity comparison.}
\label{tab:sign_control}
\begin{tabular}{lcc}
\toprule
Mode & Sign Variance & Magnitude CV \\
\midrule
Step-independent ($h=2$) & 0.18 & 0.45 \\
Episode-level & 0.00 & 0.00 \\
\bottomrule
\end{tabular}
\end{table}

%==============================================================================
\section{Method: Episode-Level GRPO}
\label{sec:method}
%==============================================================================

Our proposed fix is minimal: treat the entire episode trajectory as a single GRPO sample. Given a group of $G$ trajectories $\{\tau_1, \ldots, \tau_G\}$ with rewards $\{R_1, \ldots, R_G\}$, compute advantages $A_i = (R_i - \mu) / \sigma$ and apply the GRPO clipped objective to all tokens in $\tau_i$ with advantage $A_i$.

This eliminates both failure mechanisms simultaneously:
\begin{enumerate}
    \item \textbf{Information loss resolved}: The policy conditions on the full trajectory history during training (no truncation to $h$ steps).
    \item \textbf{Sign misattribution resolved}: All tokens share a single advantage determined by the trajectory's relative quality within the group. There is no decomposition into windows that could create conflicting gradient signals.
\end{enumerate}

The tradeoff is that episode-level credit is coarse: all actions in a trajectory are equally reinforced or suppressed regardless of their individual contribution. This limits the ability to learn fine-grained action-level strategies, as evidenced by the performance gap with GAE methods (Section~\ref{sec:gae}). We position episode-level GRPO as the \emph{necessary minimum} that eliminates the catastrophic failure, while per-token methods like GAE provide \emph{sufficient improvement} for higher performance.

%==============================================================================
\section{Experiments}
\label{sec:experiments}
%==============================================================================

\subsection{Setup}

We train Qwen2.5-1.5B-Instruct on ALFWorld using the RAGEN framework. ALFWorld provides six task types (pick, clean, heat, cool, examine, pick\_two) with average episode lengths of 18--40 steps. Training uses GRPO with group size 8, learning rate $5 \times 10^{-6}$, and binary task completion reward. All experiments run on 2$\times$A6000 GPUs.

\subsection{Episode-Level Convergence}
\label{sec:convergence}

Table~\ref{tab:convergence} shows the training trajectory of episode-level GRPO. The model converges to 37.5\% validation success at step 30, with plateau onset at step 30 (no further improvement through step 50). Pass@8 of 94\% indicates that the policy covers correct strategies but greedy sampling selection is imprecise.

\begin{table}[h]
\centering
\caption{Episode-level GRPO training trajectory (single seed). Plateau onset at step 30.}
\label{tab:convergence}
\begin{tabular}{cccccc}
\toprule
Step & Val Success & Train Success & Pass@8 & Action Valid & Entropy \\
\midrule
5 & 25.0\% & 31.3\% & 87.5\% & 0.50 & 0.22 \\
10 & 25.0\% & 37.5\% & 81.3\% & 0.52 & 0.20 \\
15 & 28.1\% & 43.8\% & 75.0\% & 0.55 & 0.19 \\
20 & 31.3\% & 50.0\% & 81.3\% & 0.57 & 0.19 \\
25 & 34.4\% & 53.1\% & 87.5\% & 0.60 & 0.18 \\
30 & 34.4\% & 56.3\% & 87.5\% & 0.61 & 0.18 \\
40 & 37.5\% & 56.3\% & 93.8\% & 0.62 & 0.18 \\
50 & 37.5\% & 56.3\% & 93.8\% & 0.62 & 0.18 \\
\bottomrule
\end{tabular}
\end{table}

\subsection{Multi-Seed Replication}

To establish reliability, we run five seeds to convergence. Table~\ref{tab:multiseed} reports results. Mean success is 34.4\% ($\pm$4.8\%), with 95\% CI [28.4\%, 40.4\%]. The lower bound of the CI (28.4\%) remains far above step-independent performance ($\leq$1\%), confirming that the episode-level advantage is robust.

\begin{table}[h]
\centering
\caption{Multi-seed episode-level GRPO results (converged). 95\% CI: [28.4\%, 40.4\%].}
\label{tab:multiseed}
\begin{tabular}{cccccc}
\toprule
Seed & Val Success & Train Success & Pass@8 & Action Valid & Converge Step \\
\midrule
1 & 37.5\% & 56.3\% & 93.8\% & 0.62 & 30 \\
2 & 31.3\% & 50.0\% & 87.5\% & 0.58 & 35 \\
3 & 28.1\% & 46.9\% & 87.5\% & 0.55 & 25 \\
4 & 40.6\% & 59.4\% & 93.8\% & 0.63 & 30 \\
5 & 34.4\% & 53.1\% & 87.5\% & 0.59 & 30 \\
\midrule
Mean$\pm$Std & 34.4$\pm$4.8\% & 53.1$\pm$4.9\% & 90.0$\pm$3.4\% & 0.59$\pm$0.04 & 30$\pm$3.5 \\
\bottomrule
\end{tabular}
\end{table}

\subsection{Context Window Sweep (h-Sweep)}
\label{sec:hsweep}

To characterize the failure boundary, we sweep context window size $h \in \{2, 5, 10, 20, \text{full}\}$ with three seeds each, training for 15 steps. Table~\ref{tab:hsweep} shows a monotonic degradation pattern rather than a sharp cliff.

\begin{table}[h]
\centering
\caption{Context window sweep results (15 steps, 3 seeds each). $h=2$ uniformly collapses; $h=5$--$10$ is a seed-dependent transition zone; $h \geq 20$ approaches episode-level performance.}
\label{tab:hsweep}
\begin{tabular}{lcccccc}
\toprule
$h$ & Seed 1 & Seed 2 & Seed 3 & Mean & Std & Collapse \\
\midrule
$h=2$ & 0.0\% & 0.0\% & 3.1\% & 1.0\% & 1.8\% & step 3--4 \\
$h=5$ & 6.3\% & 3.1\% & 12.5\% & 7.3\% & 4.7\% & step 5--7 \\
$h=10$ & 15.6\% & 12.5\% & 21.9\% & 16.7\% & 4.7\% & none \\
$h=20$ & 25.0\% & 21.9\% & 28.1\% & 25.0\% & 3.1\% & none \\
full & 31.3\% & 25.0\% & 28.1\% & 28.1\% & 3.1\% & none \\
\bottomrule
\end{tabular}
\end{table}

Key observations: (1)~$h=2$ produces robust collapse (3/3 seeds $\approx 0\%$), confirming that minimal context is catastrophic. (2)~The transition zone at $h=5$--$10$ shows high seed variance, consistent with the stochastic nature of whether a 5--10 step window can capture sufficient strategy information for a given initialization. (3)~Even $h=20$ underperforms full episode-level (25.0\% vs 28.1\%), suggesting that sign misattribution persists whenever episode reward must be decomposed across multiple windows.

Table~\ref{tab:hsweep_validity} further shows that action validity degrades monotonically with smaller $h$, indicating that both information loss (inability to select valid actions) and sign misattribution (reinforcing invalid actions from successful episodes) contribute.

\begin{table}[h]
\centering
\caption{Action validity at step 15 across context window sizes.}
\label{tab:hsweep_validity}
\begin{tabular}{lcc}
\toprule
$h$ & Mean Action Valid & Mean Invalid \\
\midrule
$h=2$ & 0.31 & 0.69 \\
$h=5$ & 0.48 & 0.52 \\
$h=10$ & 0.55 & 0.45 \\
$h=20$ & 0.58 & 0.42 \\
full & 0.60 & 0.40 \\
\bottomrule
\end{tabular}
\end{table}

\subsection{GAE Baseline Comparison}
\label{sec:gae}

To understand the performance ceiling from finer credit assignment, we compare episode-level GRPO with GAE variants that provide per-token advantages via a learned critic. Table~\ref{tab:gae} reports results at step 15 (3 seeds each).

\begin{table}[h]
\centering
\caption{Credit assignment method comparison (step 15, 3 seeds). GAE($\lambda=0.95$) provides an additional 12.5pp over episode-level GRPO, demonstrating the value of per-token credit beyond the necessary minimum.}
\label{tab:gae}
\begin{tabular}{lcccc}
\toprule
Method & Mean Success & Std & Adv. Variance & Cost Ratio \\
\midrule
Step-independent ($h=2$) & 1.0\% & 1.8\% & 0.85 & 1.0$\times$ \\
Episode-level GRPO & 28.1\% & 3.1\% & 0.00 & 1.2$\times$ \\
GAE($\lambda=1.0, \gamma=1.0$) & 34.4\% & 3.1\% & 0.28 & 1.8$\times$ \\
GAE($\lambda=0.95, \gamma=1.0$) & 40.6\% & 3.1\% & 0.42 & 1.8$\times$ \\
\bottomrule
\end{tabular}
\end{table}

The progression reveals a clean decomposition: \textbf{episode-level GRPO} eliminates sign misattribution (lifting from 1\% to 28.1\%), while \textbf{GAE} adds per-token credit that enables action-level learning (further lifting to 34.4--40.6\%). GAE($\lambda=1.0$) is equivalent to Monte Carlo returns with a baseline, providing moderate variance reduction. GAE($\lambda=0.95$) adds bootstrapping through the critic, achieving the highest performance at 1.8$\times$ training cost (due to the critic head).

This supports our framing: episode-level credit is the \emph{necessary minimum} to eliminate catastrophic failure; GAE-style per-token credit is a \emph{sufficient improvement} for higher performance. The two address orthogonal aspects of the credit assignment problem.

\subsection{Per-Task-Type Analysis}
\label{sec:pertask}

ALFWorld contains six task types with varying episode lengths. Table~\ref{tab:pertask} breaks down performance by task type, revealing that step-independent failure severity correlates strongly with episode length ($r = -0.82$).

\begin{table}[h]
\centering
\caption{Per-task-type breakdown. Step-independent failure is most severe for longer tasks ($r=-0.82$ between average steps and episode-level advantage magnitude).}
\label{tab:pertask}
\begin{tabular}{lcccccc}
\toprule
Task Type & $n$ & Avg Steps & Episode-Level & Step-Indep. & $\Delta$ \\
\midrule
pick & 40 & 18 & 50.0\% & 5.0\% & +45.0pp \\
clean & 25 & 28 & 36.0\% & 0.0\% & +36.0pp \\
cool & 15 & 25 & 33.3\% & 0.0\% & +33.3pp \\
heat & 20 & 32 & 30.0\% & 0.0\% & +30.0pp \\
pick\_two & 18 & 40 & 22.2\% & 0.0\% & +22.2pp \\
examine & 10 & 35 & 20.0\% & 0.0\% & +20.0pp \\
\bottomrule
\end{tabular}
\end{table}

The residual 5\% success for ``pick'' under step-independent training occurs because some pick tasks complete within 1--2 windows, partially avoiding sign misattribution. For all other task types requiring 25+ steps, step-independent training produces zero success. This is consistent with the mechanism: more windows mean more misattributed gradient signals, creating a stronger destructive interference pattern that prevents any learning.

\subsection{Verification Ablation}

We examine the role of self-verification in the episode-level setting. Table~\ref{tab:verify_ablation} compares always-on verification, no verification, and episode-level with an auxiliary valid\_action\_bonus.

\begin{table}[h]
\centering
\caption{Verification and reward shaping ablations. Self-verification provides +12.5pp; valid\_action\_bonus improves action validity without gaming but does not break the performance plateau.}
\label{tab:verify_ablation}
\begin{tabular}{lccccc}
\toprule
Configuration & Val Success & Action Valid & Gaming & Bonus Ratio \\
\midrule
Episode-level (no bonus) & 28.1\% & 54.5\% & N/A & N/A \\
Episode-level + bonus=0.05 & 31.3\% & 72.0\% & No & 0.22 \\
Step-independent + bonus=0.05 & 0.0\% & 97.5\% & Yes & 67.0 \\
\midrule
Always-on verification & 37.5\% & --- & --- & --- \\
No verification & 25.0\% & 52\% & --- & --- \\
\bottomrule
\end{tabular}
\end{table}

Two findings emerge. First, self-verification contributes 12.5pp (37.5\% vs 25.0\%), confirming its value for multi-turn agents. Second, a dense valid\_action\_bonus is safely absorbed under episode-level training (no gaming, bonus-to-task ratio 0.22) but catastrophically gamed under step-independent training (97.5\% action validity at 0\% success, ratio 67.0). This provides additional evidence that step-independent optimization creates pathological reward-seeking at the window level.

%==============================================================================
\section{Preliminary Results: Selective Verification}
\label{sec:selective}
%==============================================================================

Building on the diagnostic findings above, we conduct preliminary experiments on training agents to selectively verify their actions. We implement a binary verify/skip decision token and train with episode-level GRPO using action-change as an intrinsic reward for verification informativeness.

\begin{table}[h]
\centering
\caption{Selective verification comparison. B-IAVR reduces verification 28\% with 3.1pp success cost; no method dominates the Pareto frontier.}
\label{tab:selective}
\begin{tabular}{lcccc}
\toprule
Method & Success & Verif. Rate & Tokens & Cost Saving \\
\midrule
Always-on & 37.5\% & 100\% & 45K & 0\% \\
B-IAVR (ours) & 34.4\% & 72\% & 38K & 15.6\% \\
Confidence-threshold & 31.3\% & 65\% & 37K & 17.8\% \\
Random $p=0.5$ & 31.3\% & 50\% & 35K & 22.2\% \\
No-verify & 25.0\% & 0\% & 28K & 37.8\% \\
\bottomrule
\end{tabular}
\end{table}

The selective verification results are mixed. B-IAVR achieves a modest 28\% verification reduction with 3.1pp success cost, but does not dominate simple baselines on the Pareto frontier. Analysis of the action-change reward signal reveals entropy confounding: partial correlation controlling for policy entropy drops to 0.08 ($p=0.12$, not significant), suggesting that action-change partially proxies uncertainty rather than verification value.

A causal intervention experiment (verification-triggered vs random action changes at matched rates) provides weak evidence of verification utility: 37.5\% vs 28.1\% success ($p=0.04$, $n=128$ per condition). The effect is real but modest, consistent with the limited headroom at 37.5\% base success.

We treat selective verification as a preliminary direction enabled by our diagnostic contribution, rather than a primary claim. The action-change reward signal requires refinement, potentially through entropy-conditioned designs or alternative informativeness measures.

%==============================================================================
\section{Discussion}
%==============================================================================

\paragraph{Why does the community default to step-independent?}
Three factors explain the prevalence of step-independent training: (1) context window limitations of earlier models made full-episode training infeasible, (2) token-level training is the default paradigm inherited from language model pretraining, and (3) framework design choices (e.g., RAGEN's default configuration) propagate through the community without systematic evaluation of their multi-turn consequences.

\paragraph{Relationship to Turn-PPO.}
Turn-PPO identifies ``uniform advantage'' (all turns receiving identical advantage) as a problem, proposing turn-level value functions. Our sign misattribution is a superset: not only are advantages uniform within an episode, but their \emph{sign} is determined by episode outcome. Turn-PPO's solution (finer credit) and ours (coarser episode-level) address different aspects. Our GAE comparison (Section~\ref{sec:gae}) shows both have merit: episode-level eliminates misattribution (necessary), turn/token-level adds useful gradient information (beneficial).

\paragraph{Relationship to GAGPO and HiPER.}
GAGPO (93.5\%) and HiPER (97.4\%) achieve high ALFWorld performance through sophisticated credit assignment architectures. Our contribution is orthogonal: we diagnose \emph{why} the naive approach fails, not how to maximize performance. These methods implicitly solve sign misattribution through their finer credit, but the mechanism has not been previously isolated and characterized.

\paragraph{Limitations.}
Our analysis is conducted on a single environment (ALFWorld) with a 1.5B model. While the per-task-type analysis (Table~\ref{tab:pertask}) provides within-environment diversity across episode lengths, and the mechanism analysis (Section~\ref{sec:empirical_verification}) identifies general principles, verification on additional environments (e.g., ScienceWorld) and larger models (7B+) remains future work due to resource constraints. We position this as a resource-constrained diagnostic study; the identified mechanisms should hold across settings where step-independent decomposition is applied to multi-turn rewards.

%==============================================================================
\section{Conclusion}
%==============================================================================

We have identified and empirically verified a dual mechanism that causes step-independent RL training to catastrophically fail for multi-turn LLM agents: policy information loss and advantage sign misattribution. Through systematic ablation (context window sweep across 5 values with 3 seeds each, 5-seed replication, GAE comparison, per-task analysis), we demonstrate that:

\begin{enumerate}
    \item Step-independent training collapses to $\leq$1\% success on ALFWorld due to advantage signs being determined by episode outcome ($r=0.87$) rather than local action quality ($r=0.14$).
    \item Episode-level GRPO (34.4$\pm$4.8\%) provides a necessary minimum fix by eliminating sign misattribution.
    \item GAE($\lambda=0.95$) adds 12.5pp through per-token credit, representing a sufficient improvement beyond the episode-level baseline.
    \item Failure severity correlates with episode length ($r=-0.82$), consistent with the mechanism prediction.
\end{enumerate}

These findings have practical implications: researchers using RAGEN-style frameworks for multi-turn agent training should verify that their credit assignment does not introduce sign misattribution. The episode-level configuration provides a simple, effective baseline requiring no architectural changes.

\end{document}
